% Fixed subalgebras and equivariant Borcherds denominators.

\providecommand{\fg}{\mathfrak{g}}
\providecommand{\fgN}{\fg^{(N)}}
\providecommand{\fgDel}{\fg_{\Delta_5}}
\providecommand{\kBKM}{\kappa_{\mathrm{BKM}}}
\providecommand{\kch}{\kappa_{\mathrm{ch}}}
\providecommand{\kcat}{\kappa_{\mathrm{cat}}}
\providecommand{\kfib}{\kappa_{\mathrm{fiber}}}

\section{Fixed subalgebras and equivariant denominators}
\label{sec:fixed-subalgebra-twined-borcherds}

\paragraph{Statement.}
Let $\fgDel$ denote the Borcherds-Kac-Moody superalgebra whose
denominator is the primitive Gritsenko-Nikulin form $\Delta_5$ of
weight $5$, with rank-$3$ hyperbolic denominator lattice
$\Lambda^{2,1}_{II}$. Let $g_N$ be a symplectic automorphism of a K3
surface. The algebra attached to the $g_N$-twined Borcherds input is
not the fixed Lie superalgebra
\[
  \fgDel^{\langle g_N\rangle}.
\]
The former is governed by equivariant supertraces in a Borcherds
superdenominator; the latter is governed by invariant superdimensions
of root spaces.

\paragraph{Fixed roots.}
Grant the strongest comparison hypothesis: $g_N$ preserves the positive
chamber and the root label $\alpha$.
For a root $\alpha$ of $\fgDel$, write the corresponding super-root
space as $V_\alpha$. If $g_N$ acts on $V_\alpha$, the root exponent in
the ordinary denominator of the fixed subalgebra is
\[
  \operatorname{sdim} V_\alpha^{\langle g_N\rangle}
  =
  \frac{1}{N}
  \sum_{j=0}^{N-1}
  \operatorname{str}(g_N^j\mid V_\alpha).
\]
After choosing the same positive chamber, the denominator has the form
\[
  D\bigl(\fgDel^{\langle g_N\rangle}\bigr)
  =
  e^\rho
  \prod_{\alpha>0}
  (1-e^{-\alpha})^{\operatorname{sdim}V_\alpha^{\langle g_N\rangle}}.
\]
It records invariant subspaces of the untwined algebra.

\paragraph{Equivariant denominator.}
The $g_N$-equivariant superdenominator of $\fgDel$ is
\[
  D_{g_N}(\fgDel)
  =
  e^\rho
  \prod_{\alpha>0}
  \operatorname{sdet}\bigl(1-g_Ne^{-\alpha}\mid V_\alpha\bigr).
\]
Its logarithm has coefficient
\[
  -\frac{1}{m}\operatorname{str}(g_N^m\mid V_\alpha)
  \quad\text{at }e^{-m\alpha}.
\]
Thus equality
\[
  D\bigl(\fgDel^{\langle g_N\rangle}\bigr)=D_{g_N}(\fgDel)
\]
forces
\[
  \operatorname{sdim}V_\alpha^{\langle g_N\rangle}
  =
  \operatorname{str}(g_N^m\mid V_\alpha)
  \qquad (m\geq 1)
\]
for every contributing root. Equivalently, every nontrivial
$g_N$-eigenspace of every $V_\alpha$ has zero signed dimension. This is
a root-by-root condition, not a consequence of symplecticity on K3.

\paragraph{Cohomological obstruction.}
For a symplectic involution $g_2$ on a K3 surface, the holomorphic
Lefschetz trace on total cohomology is
\[
  \operatorname{Tr}(g_2\mid H^*(K3))=8
\]
and $\dim H^*(K3)=24$. Character projection gives
\[
  \dim H^*(K3)^{\langle g_2\rangle}
  =
  \frac{24+8}{2}=16.
\]
The same first K3 block therefore gives trace $8$ in the equivariant
column and invariant dimension $16$ in the fixed-subspace column. This
single computation already separates the twined denominator from the
ordinary denominator of $\fgDel^{\langle g_2\rangle}$.

\paragraph{Cartan lattice.}
The lattice $\Lambda^{2,1}_{II}$ is part of the denominator datum of
$\fgDel$. A symplectic K3 automorphism acts canonically on K3
cohomology and Mukai-lattice data. An action on $\Lambda^{2,1}_{II}$
requires an additional integral isometry preserving the positive
chamber, the Borcherds cocycle, and the denominator orientation. If the
action on $\Lambda^{2,1}_{II}$ is trivial, the fixed lattice remains
$\Lambda^{2,1}_{II}$ and produces no level-$N$ Cartan lattice. If a
nontrivial action is imposed, that action is part of the theorem. The
level-$N$ Cartan lattice of $\fgN$ is therefore separate Borcherds-lift
data, not the fixed sublattice of $\Lambda^{2,1}_{II}$.

\paragraph{Three numerical columns.}
The fixed-subalgebra, ordinary twining, and CHL denominator constants
are distinct.
\[
\begin{array}{c|c|c|c}
\text{object} & \text{formula} & N=2\text{ entry} & \text{meaning}\\
\hline
\text{fixed subalgebra}
&
\operatorname{sdim}V_\alpha^{\langle g_N\rangle}
&
16\text{ on }H^*(K3)
&
\text{invariant dimension}
\\
\text{ordinary Mathieu twining}
&
\operatorname{Tr}(g_N\mid V_\alpha)
&
8\text{ on }H^*(K3)
&
\text{trace coefficient}
\\
\text{primitive CHL denominator}
&
\kBKM(\Phi_N)=c_N(0)/2
&
c_2(0)=8,\ \kBKM(\Phi_2)=4
&
\text{Borcherds weight}
\end{array}
\]
The equality between the last two displayed $8$'s is a first-coefficient
trace normalization. It is not a fixed-subalgebra exponent and supplies
no isomorphism $\fgN\simeq\fgDel^{\langle g_N\rangle}$.

\paragraph{Primitive CHL constants.}
On the five primitive CHL frame shapes
\[
  N\in\{1,2,3,4,6\},
\]
the Vol III normalization is
\[
  (c_1(0),c_2(0),c_3(0),c_4(0),c_6(0))
  =
  (10,8,6,4,2).
\]
Borcherds' weight theorem gives
\[
  \bigl(\kBKM(\Phi_1),\kBKM(\Phi_2),\kBKM(\Phi_3),
  \kBKM(\Phi_4),\kBKM(\Phi_6)\bigr)
  =
  (5,4,3,2,1).
\]
These are automorphic weights of primitive denominator products. They
are independent of fixed-subalgebra dimensions, ordinary Mathieu
character values, $\kcat$, compact $\kch$, and $\kfib$.

\paragraph{K3xE invariant split.}
For the compact product $X=K3\times E$,
\[
  \kcat(X)=\chi(\mathcal O_X)=
  \chi(\mathcal O_{K3})\chi(\mathcal O_E)=2\cdot 0=0,
\]
and the compact Hodge-supertrace invariant is
\[
  \kch(X)=\sum_{q=0}^{3}(-1)^q h^{0,q}(X)
  =
  1-1+1-1=0.
\]
The Heisenberg-Mukai Stage-$2$ reading is separate:
\[
  \kch^{\mathrm{Heis}}(K3\times E)=3.
\]
The primitive Borcherds denominator has
\[
  \kBKM(\Delta_5)=\mathrm{wt}(\Delta_5)=c_1(0)/2=10/2=5,
\]
and the fibre lattice value is
\[
  \kfib(K3)=\operatorname{rk}H^*(K3,\mathbb Z)=24.
\]
Thus the compact $K3\times E$ package is
\[
  \{\kcat,\ \kch^{\mathrm{Heis}},\ \kBKM,\ \kfib\}
  =
  \{0,\ 3,\ 5,\ 24\},
\]
while compact $\kch(K3\times E)=0$ remains on the Hodge-supertrace
axis.

\paragraph{Diagonal-divisor catalogue.}
The Gritsenko-Clery diagonal-divisor catalogue is an independent
eight-row paramodular family. In the local Vol III normalization its
twice-weight tuple is
\[
  (10,4,6,2,4,1,3,2),
\]
corresponding to weights
\[
  (5,2,3,1,2,\tfrac12,\tfrac32,1).
\]
It shares only the level-one row $F_{1,1}=\Delta_5$ with the primitive
CHL ladder. A comparison with the CHL denominator products requires an
explicit row map and does not follow from the fixed-subalgebra
construction.

\paragraph{Conclusion.}
The assertion
\[
  \fgN \simeq \fgDel^{\langle g_N\rangle}
\]
is false for the natural cohomological action. The correct object is
the equivariant superdenominator $D_{g_N}(\fgDel)$, together with the
Borcherds product $\Phi_N$ built from the $g_N$-twined input. A theorem
identifying any fixed subalgebra with any twined denominator algebra
must supply an action on $\Lambda^{2,1}_{II}$, a compatible action on
all super-root spaces, equality between fixed exponents and Borcherds
input coefficients, and compatibility with the Borcherds cocycle and
Weyl chamber. The K3 involution computation above violates the exponent
condition in the standard cohomological model.

\paragraph{Sources.}
Borcherds 1998, Inventiones 132, Theorem 13.3, gives the weight formula
$\mathrm{wt}(\Phi)=c(0)/2$. Gritsenko 1999 and Gritsenko and Nikulin
1995 construct the CHL paramodular denominator ladder. Mukai 1988 and
Nikulin 1980 give the symplectic K3 automorphism data. Vol III
Borcherds-weight theorem and its closed-form coefficient lemma record
the five-row CHL constants. The CY-to-chiral chapter and the local
spectrum reconciliation computation record the compact $K3\times E$
invariant split.
